**Previous Research Topic:** Resilient Election & Impeachment Policy (REIP) for Trust-Governed Multi-Agent Exploration

I undertook a multi-phase research project to keep multi-agent exploration teams productive when their leader fails, hallucinates, or is compromised. Classic leader–follower swarms assume a single, permanent leader with perfect sensors. In practice—disaster zones, RF-limited tunnels, or adversarial simulations—the leader may become unreliable. When that happens, the entire team inherits bad commands, stalls, or collides. My idea was to treat leadership as a revocable privilege governed by collective trust, so the team can impeach a failing leader and elect a new one before the mission collapses.

My work had four major components:

1. **Threat model and architecture.**  
   I formalized the mission environment (GridWorld exploration with partial maps), the agents’ capabilities (limited communication radius, local sensing), and the adversarial conditions (leader hallucination, command loss). On top of the exploration controller I designed the Resilient Election & Impeachment Policy (REIP): each agent maintains a trust score for the current leader, updates it using the mismatch between predicted and observed coverage gains (CUSUM-style residuals), and flags impeachment if trust drops below a threshold. Elections select the highest-trust agent within the connected component so governance remains decentralized. I also specified fallback autonomy rules: if an agent loses commands or receives a “stay” order, it temporarily explores the nearest frontier on its own, then rejoins once the leader stabilizes.

2. **Implementation restoration and extensions.**  
   I inherited a partially complete codebase and restored critical features: deterministic spawn seeds, memoized frontier assignment, visualization throttles, trust logging, and coverage bookkeeping. I added configuration files for four benchmark modes (baseline, baseline+fallback, REIP governance-only, REIP full) under both clean and faulted environments. I reworked `scripts/compare_reip_vs_baseline.py` to accept command-line arguments, parallelize across 20 CPU cores, auto-index plot outputs, and export CSVs.

3. **Benchmarking pipeline.**  
   To make the study reproducible, I built `run_modes_experiment.py`. It now:
   - Loads the canonical configs for each mode/environment combination.
   - Runs 200 shared seeds per condition (clean/faulted) using multiprocessing with deterministic seeding.
   - Records final coverage, coverage histories, election counts, success rate to 95% coverage, and time-to-95 (steps).
   - Auto-generates bar charts (coverage, adversarial resilience), pairwise Δµ ± CI, and text summaries.
   - Stores every run with unique CSV/plot directories so older datasets remain intact.

   I ran the full suite multiple times on a 24-core workstation. This produced ~1,600 CSV rows per run plus supporting plots. I then analyzed the data to compute means, SEM, resilience (adversarial − clean), and time-to-95 success rates.

4. **Results and deliverables.**  
   - **Clean parity:** REIP clean matched baseline within 0.8 percentage points (92.5% vs 91.8%) despite running governance overhead; fallback is disabled in clean mode to keep the comparison fair.
   - **Faulted advantage:** Under hallucination + command loss, REIP governance-only and REIP full achieved ~99.1% coverage versus 88.8% for baseline, a +10.2 pp (≈11%) improvement. Resilience plots show baseline loses 3.7 pp when faults are introduced, while REIP actually gains +7.3 pp because elections remove bad leaders quickly.
   - **Governance insights:** The ablation revealed governance (trust + elections) was the dominant resilience factor in this fault model; fallback became “insurance” for heavier comm losses. This motivated a future-work section proposing harsher fault regimes to highlight fallback’s contribution.
   - **Dissemination:** I drafted an IEEE-format paper (currently under embargo for Washington State Science Fair 2025), prepared a portfolio abstract, and built high-quality figures illustrating clean vs. adversarial coverage, resilience, and governance activity.

**My contributions:** I defined the research question, restored and extended the codebase, wrote the benchmarking scripts, ran the large-scale experiments, produced the plots/statistics, and documented the findings. There was no lab sponsor—I self-directed the entire project, including version control, command-line tooling, and data analysis in Python/NumPy/Matplotlib. Detailed implementation remains proprietary until competition judging is complete, but the methodological summary and results are ready for review.

This project taught me how to combine systems thinking, software engineering, and statistical benchmarking to prove that trust-governed leadership can stabilize multi-agent exploration without requiring perfect communications or centralized oversight.